%%%%%%%%%%%%%%%%%%%%%%%%%%%%%%%%%%%%%%%%%%%%%%%%%%%%%%%%%%%%
% 14.15 COMMON MATHEMATICAL THEOREMS
% The Composition of Advanced Mathematics
%%%%%%%%%%%%%%%%%%%%%%%%%%%%%%%%%%%%%%%%%%%%%%%%%%%%%%%%%%%%

\section{Common Mathematical Theorems}
\label{sec:common-mathematical-theorems}

\prerequisite{
Mathematical notation, logic, proof methods, sets, functions,
algebra, geometry, analysis, probability, and basic topology.
}

\learningobjectives{
By the end of this reference section, the reader should be able to:
\begin{itemize}
    \item recognize frequently cited mathematical theorems by name;
    \item recall the central conclusion of each theorem;
    \item identify typical hypotheses required before applying a theorem;
    \item distinguish algebraic, analytic, geometric, topological,
          probabilistic, and computational results;
    \item connect major theorems with the chapters in which they are developed;
    \item avoid applying named theorems outside their hypotheses.
\end{itemize}
}

%%%%%%%%%%%%%%%%%%%%%%%%%%%%%%%%%%%%%%%%%%%%%%%%%%%%%%%%%%%%
\subsection{Purpose of This Reference}
\label{subsec:common-theorem-purpose}
%%%%%%%%%%%%%%%%%%%%%%%%%%%%%%%%%%%%%%%%%%%%%%%%%%%%%%%%%%%%

This section provides a compact reference to major mathematical theorems used
throughout the textbook.

It is not intended to replace the complete statements, proofs, hypotheses,
or applications developed in the corresponding chapters.

A theorem should never be applied solely from its name.

Always verify its assumptions.

%%%%%%%%%%%%%%%%%%%%%%%%%%%%%%%%%%%%%%%%%%%%%%%%%%%%%%%%%%%%
\subsection{Pythagorean Theorem}
\label{subsec:pythagorean-theorem-summary}
%%%%%%%%%%%%%%%%%%%%%%%%%%%%%%%%%%%%%%%%%%%%%%%%%%%%%%%%%%%%

For a right triangle with legs \(a,b\) and hypotenuse \(c\),

\[
\boxed{
a^2+b^2=c^2.
}
\]

The converse also holds in Euclidean geometry:

if

\[
a^2+b^2=c^2,
\]

then the angle opposite \(c\) is a right angle.

%%%%%%%%%%%%%%%%%%%%%%%%%%%%%%%%%%%%%%%%%%%%%%%%%%%%%%%%%%%%
\subsection{Fundamental Theorem of Arithmetic}
\label{subsec:fundamental-theorem-arithmetic-summary}
%%%%%%%%%%%%%%%%%%%%%%%%%%%%%%%%%%%%%%%%%%%%%%%%%%%%%%%%%%%%

Every integer

\[
n>1
\]

can be written as a product of prime numbers.

This factorization is unique up to the order of the factors:

\[
\boxed{
n
=
p_1^{\alpha_1}
\cdots
p_k^{\alpha_k}.
}
\]

%%%%%%%%%%%%%%%%%%%%%%%%%%%%%%%%%%%%%%%%%%%%%%%%%%%%%%%%%%%%
\subsection{Euclidean Algorithm}
\label{subsec:euclidean-algorithm-summary}
%%%%%%%%%%%%%%%%%%%%%%%%%%%%%%%%%%%%%%%%%%%%%%%%%%%%%%%%%%%%

For integers \(a,b\) with \(b\neq0\),

\[
\boxed{
\gcd(a,b)
=
\gcd(b,a\bmod b).
}
\]

Repeated division eventually produces the greatest common divisor.

%%%%%%%%%%%%%%%%%%%%%%%%%%%%%%%%%%%%%%%%%%%%%%%%%%%%%%%%%%%%
\subsection{Bézout's Identity}
\label{subsec:bezout-identity-summary}
%%%%%%%%%%%%%%%%%%%%%%%%%%%%%%%%%%%%%%%%%%%%%%%%%%%%%%%%%%%%

For integers \(a,b\), there exist integers \(x,y\) such that

\[
\boxed{
ax+by
=
\gcd(a,b).
}
\]

In particular,

\[
\gcd(a,b)=1
\]

if and only if there exist integers \(x,y\) with

\[
\boxed{
ax+by=1.
}
\]

%%%%%%%%%%%%%%%%%%%%%%%%%%%%%%%%%%%%%%%%%%%%%%%%%%%%%%%%%%%%
\subsection{Euclid's Lemma}
\label{subsec:euclids-lemma-summary}
%%%%%%%%%%%%%%%%%%%%%%%%%%%%%%%%%%%%%%%%%%%%%%%%%%%%%%%%%%%%

If \(p\) is prime and

\[
p\mid ab,
\]

then

\[
\boxed{
p\mid a
\quad
\text{or}
\quad
p\mid b.
}
\]

%%%%%%%%%%%%%%%%%%%%%%%%%%%%%%%%%%%%%%%%%%%%%%%%%%%%%%%%%%%%
\subsection{Fermat's Little Theorem}
\label{subsec:fermats-little-theorem-summary}
%%%%%%%%%%%%%%%%%%%%%%%%%%%%%%%%%%%%%%%%%%%%%%%%%%%%%%%%%%%%

If \(p\) is prime and

\[
p\nmid a,
\]

then

\[
\boxed{
a^{p-1}
\equiv
1
\pmod p.
}
\]

Equivalent form:

\[
\boxed{
a^p
\equiv
a
\pmod p
}
\]

for every integer \(a\).

%%%%%%%%%%%%%%%%%%%%%%%%%%%%%%%%%%%%%%%%%%%%%%%%%%%%%%%%%%%%
\subsection{Euler's Theorem}
\label{subsec:eulers-theorem-number-summary}
%%%%%%%%%%%%%%%%%%%%%%%%%%%%%%%%%%%%%%%%%%%%%%%%%%%%%%%%%%%%

If

\[
\gcd(a,n)=1,
\]

then

\[
\boxed{
a^{\varphi(n)}
\equiv
1
\pmod n,
}
\]

where \(\varphi\) is Euler's totient function.

%%%%%%%%%%%%%%%%%%%%%%%%%%%%%%%%%%%%%%%%%%%%%%%%%%%%%%%%%%%%
\subsection{Chinese Remainder Theorem}
\label{subsec:crt-summary}
%%%%%%%%%%%%%%%%%%%%%%%%%%%%%%%%%%%%%%%%%%%%%%%%%%%%%%%%%%%%

If

\[
n_1,\ldots,n_k
\]

are pairwise coprime, then the simultaneous congruences

\[
x\equiv a_i\pmod{n_i}
\]

have a unique solution modulo

\[
\boxed{
N=n_1n_2\cdots n_k.
}
\]

%%%%%%%%%%%%%%%%%%%%%%%%%%%%%%%%%%%%%%%%%%%%%%%%%%%%%%%%%%%%
\subsection{Fundamental Theorem of Algebra}
\label{subsec:fundamental-theorem-algebra-summary}
%%%%%%%%%%%%%%%%%%%%%%%%%%%%%%%%%%%%%%%%%%%%%%%%%%%%%%%%%%%%

Every nonconstant polynomial with complex coefficients has at least one
complex root.

Equivalently, every degree-\(n\) polynomial over \(\mathbb{C}\) factors as

\[
\boxed{
p(z)
=
a
(z-z_1)\cdots(z-z_n)
}
\]

counting roots with multiplicity.

%%%%%%%%%%%%%%%%%%%%%%%%%%%%%%%%%%%%%%%%%%%%%%%%%%%%%%%%%%%%
\subsection{Factor Theorem}
\label{subsec:factor-theorem-summary}
%%%%%%%%%%%%%%%%%%%%%%%%%%%%%%%%%%%%%%%%%%%%%%%%%%%%%%%%%%%%

For a polynomial \(p(x)\),

\[
\boxed{
p(a)=0
}
\]

if and only if

\[
\boxed{
x-a
}
\]

is a factor of \(p(x)\).

%%%%%%%%%%%%%%%%%%%%%%%%%%%%%%%%%%%%%%%%%%%%%%%%%%%%%%%%%%%%
\subsection{Remainder Theorem}
\label{subsec:remainder-theorem-summary}
%%%%%%%%%%%%%%%%%%%%%%%%%%%%%%%%%%%%%%%%%%%%%%%%%%%%%%%%%%%%

The remainder when polynomial \(p(x)\) is divided by \(x-a\) is

\[
\boxed{
p(a).
}
\]

%%%%%%%%%%%%%%%%%%%%%%%%%%%%%%%%%%%%%%%%%%%%%%%%%%%%%%%%%%%%
\subsection{Binomial Theorem}
\label{subsec:binomial-theorem-summary}
%%%%%%%%%%%%%%%%%%%%%%%%%%%%%%%%%%%%%%%%%%%%%%%%%%%%%%%%%%%%

For nonnegative integer \(n\),

\[
\boxed{
(a+b)^n
=
\sum_{k=0}^{n}
\binom{n}{k}
a^{n-k}b^k.
}
\]

%%%%%%%%%%%%%%%%%%%%%%%%%%%%%%%%%%%%%%%%%%%%%%%%%%%%%%%%%%%%
\subsection{Rank--Nullity Theorem}
\label{subsec:rank-nullity-summary}
%%%%%%%%%%%%%%%%%%%%%%%%%%%%%%%%%%%%%%%%%%%%%%%%%%%%%%%%%%%%

For a linear map

\[
T:V\to W
\]

with finite-dimensional \(V\),

\[
\boxed{
\dim V
=
\operatorname{rank}(T)
+
\operatorname{nullity}(T).
}
\]

Equivalently,

\[
\boxed{
\dim V
=
\dim(\operatorname{im}T)
+
\dim(\ker T).
}
\]

%%%%%%%%%%%%%%%%%%%%%%%%%%%%%%%%%%%%%%%%%%%%%%%%%%%%%%%%%%%%
\subsection{Invertible Matrix Theorem}
\label{subsec:invertible-matrix-summary}
%%%%%%%%%%%%%%%%%%%%%%%%%%%%%%%%%%%%%%%%%%%%%%%%%%%%%%%%%%%%

For an \(n\times n\) matrix \(A\), many conditions are equivalent, including:

\[
\boxed{
A
\text{ is invertible},
}
\]

\[
\boxed{
\det A\neq0,
}
\]

\[
\boxed{
\operatorname{rank}A=n,
}
\]

\[
\boxed{
\ker A=\{\mathbf0\},
}
\]

and

\[
\boxed{
0
\text{ is not an eigenvalue of }A.
}
\]

%%%%%%%%%%%%%%%%%%%%%%%%%%%%%%%%%%%%%%%%%%%%%%%%%%%%%%%%%%%%
\subsection{Cayley--Hamilton Theorem}
\label{subsec:cayley-hamilton-summary}
%%%%%%%%%%%%%%%%%%%%%%%%%%%%%%%%%%%%%%%%%%%%%%%%%%%%%%%%%%%%

Every square matrix satisfies its own characteristic polynomial.

If

\[
p_A(\lambda)
=
\det(\lambda I-A),
\]

then

\[
\boxed{
p_A(A)=0.
}
\]

%%%%%%%%%%%%%%%%%%%%%%%%%%%%%%%%%%%%%%%%%%%%%%%%%%%%%%%%%%%%
\subsection{Spectral Theorem for Real Symmetric Matrices}
\label{subsec:spectral-theorem-symmetric-summary}
%%%%%%%%%%%%%%%%%%%%%%%%%%%%%%%%%%%%%%%%%%%%%%%%%%%%%%%%%%%%

Every real symmetric matrix \(A\) has:

\[
\boxed{
\text{real eigenvalues},
}
\]

and an orthonormal basis of eigenvectors.

Thus

\[
\boxed{
A
=
QDQ^T
}
\]

for orthogonal \(Q\) and diagonal \(D\).

%%%%%%%%%%%%%%%%%%%%%%%%%%%%%%%%%%%%%%%%%%%%%%%%%%%%%%%%%%%%
\subsection{Spectral Theorem for Normal Operators}
\label{subsec:spectral-theorem-normal-summary}
%%%%%%%%%%%%%%%%%%%%%%%%%%%%%%%%%%%%%%%%%%%%%%%%%%%%%%%%%%%%

A complex finite-dimensional normal matrix is unitarily diagonalizable.

If

\[
AA^*=A^*A,
\]

then

\[
\boxed{
A
=
UDU^*
}
\]

for unitary \(U\) and diagonal \(D\).

%%%%%%%%%%%%%%%%%%%%%%%%%%%%%%%%%%%%%%%%%%%%%%%%%%%%%%%%%%%%
\subsection{Singular Value Decomposition}
\label{subsec:svd-theorem-summary}
%%%%%%%%%%%%%%%%%%%%%%%%%%%%%%%%%%%%%%%%%%%%%%%%%%%%%%%%%%%%

Every real \(m\times n\) matrix \(A\) can be written

\[
\boxed{
A
=
U\Sigma V^T,
}
\]

where \(U,V\) are orthogonal and \(\Sigma\) is diagonal in the rectangular
sense with nonnegative singular values.

%%%%%%%%%%%%%%%%%%%%%%%%%%%%%%%%%%%%%%%%%%%%%%%%%%%%%%%%%%%%
\subsection{Fundamental Theorem of Homomorphisms}
\label{subsec:first-isomorphism-summary}
%%%%%%%%%%%%%%%%%%%%%%%%%%%%%%%%%%%%%%%%%%%%%%%%%%%%%%%%%%%%

For a homomorphism

\[
\varphi:G\to H,
\]

\[
\boxed{
G/\ker\varphi
\cong
\operatorname{im}\varphi.
}
\]

Analogous first-isomorphism theorems hold for rings, modules, and other
algebraic structures.

%%%%%%%%%%%%%%%%%%%%%%%%%%%%%%%%%%%%%%%%%%%%%%%%%%%%%%%%%%%%
\subsection{Lagrange's Theorem}
\label{subsec:lagranges-theorem-group-summary}
%%%%%%%%%%%%%%%%%%%%%%%%%%%%%%%%%%%%%%%%%%%%%%%%%%%%%%%%%%%%

If \(G\) is a finite group and \(H\leq G\), then

\[
\boxed{
|H|
\mid
|G|.
}
\]

Moreover,

\[
\boxed{
|G|
=
[G:H]|H|.
}
\]

%%%%%%%%%%%%%%%%%%%%%%%%%%%%%%%%%%%%%%%%%%%%%%%%%%%%%%%%%%%%
\subsection{Cauchy's Theorem for Finite Groups}
\label{subsec:cauchy-group-summary}
%%%%%%%%%%%%%%%%%%%%%%%%%%%%%%%%%%%%%%%%%%%%%%%%%%%%%%%%%%%%

If prime \(p\) divides \(|G|\), then \(G\) contains an element of order \(p\).

%%%%%%%%%%%%%%%%%%%%%%%%%%%%%%%%%%%%%%%%%%%%%%%%%%%%%%%%%%%%
\subsection{Sylow Theorems}
\label{subsec:sylow-summary}
%%%%%%%%%%%%%%%%%%%%%%%%%%%%%%%%%%%%%%%%%%%%%%%%%%%%%%%%%%%%

If

\[
|G|
=
p^nm,
\qquad
p\nmid m,
\]

then finite group \(G\) has subgroups of order \(p^n\), called Sylow
\(p\)-subgroups.

The Sylow theorems also describe their conjugacy and constrain their number.

%%%%%%%%%%%%%%%%%%%%%%%%%%%%%%%%%%%%%%%%%%%%%%%%%%%%%%%%%%%%
\subsection{Orbit--Stabilizer Theorem}
\label{subsec:orbit-stabilizer-summary}
%%%%%%%%%%%%%%%%%%%%%%%%%%%%%%%%%%%%%%%%%%%%%%%%%%%%%%%%%%%%

If group \(G\) acts on a set and \(x\) is an element, then

\[
\boxed{
|G\cdot x|
=
[G:G_x],
}
\]

where \(G_x\) is the stabilizer of \(x\).

For finite \(G\),

\[
\boxed{
|G|
=
|G\cdot x|
\,|G_x|.
}
\]

%%%%%%%%%%%%%%%%%%%%%%%%%%%%%%%%%%%%%%%%%%%%%%%%%%%%%%%%%%%%
\subsection{Law of Sines}
\label{subsec:law-sines-theorem-summary}
%%%%%%%%%%%%%%%%%%%%%%%%%%%%%%%%%%%%%%%%%%%%%%%%%%%%%%%%%%%%

For a Euclidean triangle,

\[
\boxed{
\frac{a}{\sin A}
=
\frac{b}{\sin B}
=
\frac{c}{\sin C}
=
2R,
}
\]

where \(R\) is the circumradius.

%%%%%%%%%%%%%%%%%%%%%%%%%%%%%%%%%%%%%%%%%%%%%%%%%%%%%%%%%%%%
\subsection{Law of Cosines}
\label{subsec:law-cosines-theorem-summary}
%%%%%%%%%%%%%%%%%%%%%%%%%%%%%%%%%%%%%%%%%%%%%%%%%%%%%%%%%%%%

For triangle sides \(a,b,c\),

\[
\boxed{
c^2
=
a^2+b^2-2ab\cos C.
}
\]

The Pythagorean theorem is the special case

\[
C=\frac{\pi}{2}.
\]

%%%%%%%%%%%%%%%%%%%%%%%%%%%%%%%%%%%%%%%%%%%%%%%%%%%%%%%%%%%%
\subsection{Ceva's Theorem}
\label{subsec:ceva-summary}
%%%%%%%%%%%%%%%%%%%%%%%%%%%%%%%%%%%%%%%%%%%%%%%%%%%%%%%%%%%%

For points \(D,E,F\) on the sides of triangle \(ABC\), the cevians
\(AD,BE,CF\) are concurrent if and only if, with appropriate directed or
positive-length conventions,

\[
\boxed{
\frac{BD}{DC}
\frac{CE}{EA}
\frac{AF}{FB}
=
1.
}
\]

%%%%%%%%%%%%%%%%%%%%%%%%%%%%%%%%%%%%%%%%%%%%%%%%%%%%%%%%%%%%
\subsection{Menelaus' Theorem}
\label{subsec:menelaus-summary}
%%%%%%%%%%%%%%%%%%%%%%%%%%%%%%%%%%%%%%%%%%%%%%%%%%%%%%%%%%%%

For a transversal intersecting the sides or extended sides of a triangle,
the resulting directed ratios satisfy a product relation.

Under a standard directed-length convention,

\[
\boxed{
\frac{AF}{FB}
\frac{BD}{DC}
\frac{CE}{EA}
=
-1.
}
\]

%%%%%%%%%%%%%%%%%%%%%%%%%%%%%%%%%%%%%%%%%%%%%%%%%%%%%%%%%%%%
\subsection{Euler Polyhedron Formula}
\label{subsec:euler-polyhedron-summary}
%%%%%%%%%%%%%%%%%%%%%%%%%%%%%%%%%%%%%%%%%%%%%%%%%%%%%%%%%%%%

For a connected convex polyhedron,

\[
\boxed{
V-E+F=2,
}
\]

where \(V,E,F\) denote numbers of vertices, edges, and faces.

More generally, Euler characteristic extends this idea to topological spaces.

%%%%%%%%%%%%%%%%%%%%%%%%%%%%%%%%%%%%%%%%%%%%%%%%%%%%%%%%%%%%
\subsection{Jordan Curve Theorem}
\label{subsec:jordan-curve-summary}
%%%%%%%%%%%%%%%%%%%%%%%%%%%%%%%%%%%%%%%%%%%%%%%%%%%%%%%%%%%%

A simple closed curve in the plane separates the plane into exactly two
components:

\[
\boxed{
\text{an interior and an exterior},
}
\]

with the curve as their common boundary.

%%%%%%%%%%%%%%%%%%%%%%%%%%%%%%%%%%%%%%%%%%%%%%%%%%%%%%%%%%%%
\subsection{Brouwer Fixed-Point Theorem}
\label{subsec:brouwer-fixed-point-summary}
%%%%%%%%%%%%%%%%%%%%%%%%%%%%%%%%%%%%%%%%%%%%%%%%%%%%%%%%%%%%

Every continuous function from a closed ball in \(\mathbb{R}^n\) to itself
has at least one fixed point.

That is, there exists \(x\) such that

\[
\boxed{
f(x)=x.
}
\]

%%%%%%%%%%%%%%%%%%%%%%%%%%%%%%%%%%%%%%%%%%%%%%%%%%%%%%%%%%%%
\subsection{Banach Fixed-Point Theorem}
\label{subsec:banach-fixed-point-summary}
%%%%%%%%%%%%%%%%%%%%%%%%%%%%%%%%%%%%%%%%%%%%%%%%%%%%%%%%%%%%

A contraction mapping on a complete metric space has a unique fixed point.

If

\[
d(Tx,Ty)
\leq
q\,d(x,y),
\qquad
0\leq q<1,
\]

then repeated iteration

\[
x_{n+1}=T(x_n)
\]

converges to the unique fixed point.

%%%%%%%%%%%%%%%%%%%%%%%%%%%%%%%%%%%%%%%%%%%%%%%%%%%%%%%%%%%%
\subsection{Heine--Borel Theorem}
\label{subsec:heine-borel-summary}
%%%%%%%%%%%%%%%%%%%%%%%%%%%%%%%%%%%%%%%%%%%%%%%%%%%%%%%%%%%%

A subset of \(\mathbb{R}^n\) is compact if and only if it is closed and
bounded.

This equivalence is special to finite-dimensional Euclidean space.

%%%%%%%%%%%%%%%%%%%%%%%%%%%%%%%%%%%%%%%%%%%%%%%%%%%%%%%%%%%%
\subsection{Bolzano--Weierstrass Theorem}
\label{subsec:bolzano-weierstrass-summary}
%%%%%%%%%%%%%%%%%%%%%%%%%%%%%%%%%%%%%%%%%%%%%%%%%%%%%%%%%%%%

Every bounded sequence in \(\mathbb{R}^n\) has a convergent subsequence.

%%%%%%%%%%%%%%%%%%%%%%%%%%%%%%%%%%%%%%%%%%%%%%%%%%%%%%%%%%%%
\subsection{Extreme Value Theorem}
\label{subsec:extreme-value-theorem-summary}
%%%%%%%%%%%%%%%%%%%%%%%%%%%%%%%%%%%%%%%%%%%%%%%%%%%%%%%%%%%%

If

\[
f:[a,b]\to\mathbb{R}
\]

is continuous, then \(f\) attains both a maximum and minimum on \([a,b]\).

There exist \(x_{\min},x_{\max}\in[a,b]\) such that

\[
\boxed{
f(x_{\min})
\leq
f(x)
\leq
f(x_{\max})
}
\]

for every \(x\in[a,b]\).

%%%%%%%%%%%%%%%%%%%%%%%%%%%%%%%%%%%%%%%%%%%%%%%%%%%%%%%%%%%%
\subsection{Intermediate Value Theorem}
\label{subsec:intermediate-value-summary}
%%%%%%%%%%%%%%%%%%%%%%%%%%%%%%%%%%%%%%%%%%%%%%%%%%%%%%%%%%%%

If \(f\) is continuous on \([a,b]\), then every value between \(f(a)\) and
\(f(b)\) is attained.

If \(L\) lies between \(f(a)\) and \(f(b)\), then there exists \(c\in[a,b]\)
such that

\[
\boxed{
f(c)=L.
}
\]

%%%%%%%%%%%%%%%%%%%%%%%%%%%%%%%%%%%%%%%%%%%%%%%%%%%%%%%%%%%%
\subsection{Bolzano's Theorem}
\label{subsec:bolzano-theorem-summary}
%%%%%%%%%%%%%%%%%%%%%%%%%%%%%%%%%%%%%%%%%%%%%%%%%%%%%%%%%%%%

If \(f\) is continuous on \([a,b]\) and

\[
f(a)f(b)<0,
\]

then there exists \(c\in(a,b)\) such that

\[
\boxed{
f(c)=0.
}
\]

This is a special case of the Intermediate Value Theorem.

%%%%%%%%%%%%%%%%%%%%%%%%%%%%%%%%%%%%%%%%%%%%%%%%%%%%%%%%%%%%
\subsection{Rolle's Theorem}
\label{subsec:rolles-theorem-summary}
%%%%%%%%%%%%%%%%%%%%%%%%%%%%%%%%%%%%%%%%%%%%%%%%%%%%%%%%%%%%

If \(f\) is continuous on \([a,b]\), differentiable on \((a,b)\), and

\[
f(a)=f(b),
\]

then there exists \(c\in(a,b)\) such that

\[
\boxed{
f'(c)=0.
}
\]

%%%%%%%%%%%%%%%%%%%%%%%%%%%%%%%%%%%%%%%%%%%%%%%%%%%%%%%%%%%%
\subsection{Mean Value Theorem}
\label{subsec:mean-value-theorem-summary}
%%%%%%%%%%%%%%%%%%%%%%%%%%%%%%%%%%%%%%%%%%%%%%%%%%%%%%%%%%%%

If \(f\) is continuous on \([a,b]\) and differentiable on \((a,b)\), then
there exists \(c\in(a,b)\) such that

\[
\boxed{
f'(c)
=
\frac{
f(b)-f(a)
}{
b-a
}.
}
\]

%%%%%%%%%%%%%%%%%%%%%%%%%%%%%%%%%%%%%%%%%%%%%%%%%%%%%%%%%%%%
\subsection{Cauchy's Mean Value Theorem}
\label{subsec:cauchy-mvt-summary}
%%%%%%%%%%%%%%%%%%%%%%%%%%%%%%%%%%%%%%%%%%%%%%%%%%%%%%%%%%%%

Under appropriate continuity and differentiability assumptions on \(f,g\),
there exists \(c\in(a,b)\) such that

\[
\boxed{
[f(b)-f(a)]g'(c)
=
[g(b)-g(a)]f'(c).
}
\]

This generalizes the ordinary Mean Value Theorem.

%%%%%%%%%%%%%%%%%%%%%%%%%%%%%%%%%%%%%%%%%%%%%%%%%%%%%%%%%%%%
\subsection{Taylor's Theorem}
\label{subsec:taylors-theorem-summary}
%%%%%%%%%%%%%%%%%%%%%%%%%%%%%%%%%%%%%%%%%%%%%%%%%%%%%%%%%%%%

Under appropriate differentiability assumptions,

\[
\boxed{
f(x)
=
\sum_{k=0}^{n}
\frac{
f^{(k)}(a)
}{
k!
}
(x-a)^k
+
R_n(x).
}
\]

In Lagrange form,

\[
\boxed{
R_n(x)
=
\frac{
f^{(n+1)}(\xi)
}{
(n+1)!
}
(x-a)^{n+1}
}
\]

for some \(\xi\) between \(a\) and \(x\).

%%%%%%%%%%%%%%%%%%%%%%%%%%%%%%%%%%%%%%%%%%%%%%%%%%%%%%%%%%%%
\subsection{Fundamental Theorem of Calculus}
\label{subsec:fundamental-theorem-calculus-summary}
%%%%%%%%%%%%%%%%%%%%%%%%%%%%%%%%%%%%%%%%%%%%%%%%%%%%%%%%%%%%

If \(f\) is continuous and

\[
F(x)
=
\int_a^x
f(t)\,dt,
\]

then

\[
\boxed{
F'(x)=f(x).
}
\]

If \(G'=f\), then

\[
\boxed{
\int_a^b
f(x)\,dx
=
G(b)-G(a).
}
\]

%%%%%%%%%%%%%%%%%%%%%%%%%%%%%%%%%%%%%%%%%%%%%%%%%%%%%%%%%%%%
\subsection{Inverse Function Theorem}
\label{subsec:inverse-function-theorem-summary}
%%%%%%%%%%%%%%%%%%%%%%%%%%%%%%%%%%%%%%%%%%%%%%%%%%%%%%%%%%%%

For a continuously differentiable map

\[
F:\mathbb{R}^n\to\mathbb{R}^n,
\]

if the Jacobian matrix \(DF(a)\) is invertible, then \(F\) has a local
differentiable inverse near \(a\).

Moreover,

\[
\boxed{
D(F^{-1})(F(a))
=
[DF(a)]^{-1}.
}
\]

%%%%%%%%%%%%%%%%%%%%%%%%%%%%%%%%%%%%%%%%%%%%%%%%%%%%%%%%%%%%
\subsection{Implicit Function Theorem}
\label{subsec:implicit-function-theorem-summary}
%%%%%%%%%%%%%%%%%%%%%%%%%%%%%%%%%%%%%%%%%%%%%%%%%%%%%%%%%%%%

Under suitable differentiability and nonsingularity conditions, an equation

\[
F(x,y)=0
\]

can locally be solved for some variables as differentiable functions of the
others.

In the two-variable scalar case, if

\[
F_y(a,b)\neq0,
\]

then locally

\[
y=g(x)
\]

and

\[
\boxed{
g'(x)
=
-\frac{
F_x
}{
F_y
}.
}
\]

%%%%%%%%%%%%%%%%%%%%%%%%%%%%%%%%%%%%%%%%%%%%%%%%%%%%%%%%%%%%
\subsection{Green's Theorem}
\label{subsec:greens-theorem-summary}
%%%%%%%%%%%%%%%%%%%%%%%%%%%%%%%%%%%%%%%%%%%%%%%%%%%%%%%%%%%%

For a positively oriented simple closed curve \(C\) bounding region \(D\),

\[
\boxed{
\oint_C
P\,dx+Q\,dy
=
\iint_D
\left(
\frac{\partial Q}{\partial x}
-
\frac{\partial P}{\partial y}
\right)
dA
}
\]

under suitable regularity assumptions.

%%%%%%%%%%%%%%%%%%%%%%%%%%%%%%%%%%%%%%%%%%%%%%%%%%%%%%%%%%%%
\subsection{Divergence Theorem}
\label{subsec:divergence-theorem-summary}
%%%%%%%%%%%%%%%%%%%%%%%%%%%%%%%%%%%%%%%%%%%%%%%%%%%%%%%%%%%%

For a suitable vector field \(\mathbf{F}\) on volume \(V\) with boundary
surface \(S\),

\[
\boxed{
\iiint_V
\nabla\cdot\mathbf{F}
\,dV
=
\iint_S
\mathbf{F}\cdot\mathbf{n}
\,dS.
}
\]

%%%%%%%%%%%%%%%%%%%%%%%%%%%%%%%%%%%%%%%%%%%%%%%%%%%%%%%%%%%%
\subsection{Stokes' Theorem}
\label{subsec:stokes-theorem-summary}
%%%%%%%%%%%%%%%%%%%%%%%%%%%%%%%%%%%%%%%%%%%%%%%%%%%%%%%%%%%%

For oriented surface \(S\) with boundary \(\partial S\),

\[
\boxed{
\int_{\partial S}
\mathbf{F}\cdot d\mathbf{r}
=
\iint_S
(\nabla\times\mathbf{F})
\cdot\mathbf{n}
\,dS.
}
\]

Green's theorem is a planar special case.

%%%%%%%%%%%%%%%%%%%%%%%%%%%%%%%%%%%%%%%%%%%%%%%%%%%%%%%%%%%%
\subsection{Fubini's Theorem}
\label{subsec:fubini-summary}
%%%%%%%%%%%%%%%%%%%%%%%%%%%%%%%%%%%%%%%%%%%%%%%%%%%%%%%%%%%%

Under suitable integrability assumptions, a multiple integral may be
evaluated as iterated integrals and the order of integration may be changed.

For example,

\[
\boxed{
\iint
f(x,y)\,dx\,dy
=
\int
\left(
\int
f(x,y)\,dx
\right)
dy.
}
\]

%%%%%%%%%%%%%%%%%%%%%%%%%%%%%%%%%%%%%%%%%%%%%%%%%%%%%%%%%%%%
\subsection{Tonelli's Theorem}
\label{subsec:tonelli-summary}
%%%%%%%%%%%%%%%%%%%%%%%%%%%%%%%%%%%%%%%%%%%%%%%%%%%%%%%%%%%%

For a nonnegative measurable function, iterated integrals can be interchanged,
possibly taking the value \(+\infty\).

Tonelli's theorem is especially useful before absolute integrability has been
established.

%%%%%%%%%%%%%%%%%%%%%%%%%%%%%%%%%%%%%%%%%%%%%%%%%%%%%%%%%%%%
\subsection{Monotone Convergence Theorem}
\label{subsec:monotone-convergence-theorem-summary}
%%%%%%%%%%%%%%%%%%%%%%%%%%%%%%%%%%%%%%%%%%%%%%%%%%%%%%%%%%%%

If

\[
0\leq f_1\leq f_2\leq\cdots
\]

and

\[
f_n\to f
\]

pointwise, then

\[
\boxed{
\int f_n
\to
\int f.
}
\]

%%%%%%%%%%%%%%%%%%%%%%%%%%%%%%%%%%%%%%%%%%%%%%%%%%%%%%%%%%%%
\subsection{Dominated Convergence Theorem}
\label{subsec:dominated-convergence-summary}
%%%%%%%%%%%%%%%%%%%%%%%%%%%%%%%%%%%%%%%%%%%%%%%%%%%%%%%%%%%%

If

\[
f_n\to f
\]

almost everywhere and

\[
|f_n|\leq g
\]

for an integrable function \(g\), then

\[
\boxed{
\int f_n
\to
\int f.
}
\]

Under these assumptions,

\[
f
\]

is also integrable.

%%%%%%%%%%%%%%%%%%%%%%%%%%%%%%%%%%%%%%%%%%%%%%%%%%%%%%%%%%%%
\subsection{Fatou's Lemma}
\label{subsec:fatou-summary}
%%%%%%%%%%%%%%%%%%%%%%%%%%%%%%%%%%%%%%%%%%%%%%%%%%%%%%%%%%%%

For nonnegative measurable \(f_n\),

\[
\boxed{
\int
\liminf_{n\to\infty}f_n
\leq
\liminf_{n\to\infty}
\int f_n.
}
\]

%%%%%%%%%%%%%%%%%%%%%%%%%%%%%%%%%%%%%%%%%%%%%%%%%%%%%%%%%%%%
\subsection{Cauchy--Schwarz Inequality}
\label{subsec:cauchy-schwarz-summary}
%%%%%%%%%%%%%%%%%%%%%%%%%%%%%%%%%%%%%%%%%%%%%%%%%%%%%%%%%%%%

For vectors in an inner-product space,

\[
\boxed{
|\langle u,v\rangle|
\leq
\|u\|\|v\|.
}
\]

Equality occurs precisely when \(u,v\) are linearly dependent, including the
zero-vector cases.

%%%%%%%%%%%%%%%%%%%%%%%%%%%%%%%%%%%%%%%%%%%%%%%%%%%%%%%%%%%%
\subsection{Triangle Inequality}
\label{subsec:triangle-inequality-summary}
%%%%%%%%%%%%%%%%%%%%%%%%%%%%%%%%%%%%%%%%%%%%%%%%%%%%%%%%%%%%

For a norm,

\[
\boxed{
\|u+v\|
\leq
\|u\|+\|v\|.
}
\]

For real or complex numbers,

\[
\boxed{
|a+b|
\leq
|a|+|b|.
}
\]

%%%%%%%%%%%%%%%%%%%%%%%%%%%%%%%%%%%%%%%%%%%%%%%%%%%%%%%%%%%%
\subsection{Hölder's Inequality}
\label{subsec:holder-summary}
%%%%%%%%%%%%%%%%%%%%%%%%%%%%%%%%%%%%%%%%%%%%%%%%%%%%%%%%%%%%

If

\[
1<p,q<\infty,
\qquad
\frac1p+\frac1q=1,
\]

then

\[
\boxed{
\int
|fg|
\leq
\left(
\int|f|^p
\right)^{1/p}
\left(
\int|g|^q
\right)^{1/q}.
}
\]

%%%%%%%%%%%%%%%%%%%%%%%%%%%%%%%%%%%%%%%%%%%%%%%%%%%%%%%%%%%%
\subsection{Minkowski's Inequality}
\label{subsec:minkowski-summary}
%%%%%%%%%%%%%%%%%%%%%%%%%%%%%%%%%%%%%%%%%%%%%%%%%%%%%%%%%%%%

For \(p\geq1\),

\[
\boxed{
\|f+g\|_p
\leq
\|f\|_p+\|g\|_p.
}
\]

This is the triangle inequality for \(L^p\) spaces.

%%%%%%%%%%%%%%%%%%%%%%%%%%%%%%%%%%%%%%%%%%%%%%%%%%%%%%%%%%%%
\subsection{Baire Category Theorem}
\label{subsec:baire-category-summary}
%%%%%%%%%%%%%%%%%%%%%%%%%%%%%%%%%%%%%%%%%%%%%%%%%%%%%%%%%%%%

A complete metric space cannot be expressed as a countable union of nowhere
dense closed sets.

This theorem underlies several major results in functional analysis.

%%%%%%%%%%%%%%%%%%%%%%%%%%%%%%%%%%%%%%%%%%%%%%%%%%%%%%%%%%%%
\subsection{Banach--Steinhaus Theorem}
\label{subsec:uniform-boundedness-summary}
%%%%%%%%%%%%%%%%%%%%%%%%%%%%%%%%%%%%%%%%%%%%%%%%%%%%%%%%%%%%

For a family of bounded linear operators on a Banach space, pointwise
boundedness implies uniform boundedness in operator norm under the standard
hypotheses.

This is also called the Uniform Boundedness Principle.

%%%%%%%%%%%%%%%%%%%%%%%%%%%%%%%%%%%%%%%%%%%%%%%%%%%%%%%%%%%%
\subsection{Open Mapping Theorem}
\label{subsec:open-mapping-summary}
%%%%%%%%%%%%%%%%%%%%%%%%%%%%%%%%%%%%%%%%%%%%%%%%%%%%%%%%%%%%

A surjective bounded linear operator between Banach spaces is an open map.

%%%%%%%%%%%%%%%%%%%%%%%%%%%%%%%%%%%%%%%%%%%%%%%%%%%%%%%%%%%%
\subsection{Closed Graph Theorem}
\label{subsec:closed-graph-summary}
%%%%%%%%%%%%%%%%%%%%%%%%%%%%%%%%%%%%%%%%%%%%%%%%%%%%%%%%%%%%

A linear operator between Banach spaces whose graph is closed is bounded,
provided the operator is defined on the entire domain space.

%%%%%%%%%%%%%%%%%%%%%%%%%%%%%%%%%%%%%%%%%%%%%%%%%%%%%%%%%%%%
\subsection{Hahn--Banach Theorem}
\label{subsec:hahn-banach-summary}
%%%%%%%%%%%%%%%%%%%%%%%%%%%%%%%%%%%%%%%%%%%%%%%%%%%%%%%%%%%%

Under standard hypotheses, a bounded linear functional defined on a subspace
can be extended to the entire vector space without increasing its norm.

%%%%%%%%%%%%%%%%%%%%%%%%%%%%%%%%%%%%%%%%%%%%%%%%%%%%%%%%%%%%
\subsection{Cauchy's Integral Theorem}
\label{subsec:cauchy-integral-theorem-summary}
%%%%%%%%%%%%%%%%%%%%%%%%%%%%%%%%%%%%%%%%%%%%%%%%%%%%%%%%%%%%

For a holomorphic function on an appropriate simply connected region,

\[
\boxed{
\oint_C
f(z)\,dz
=
0
}
\]

for closed contours \(C\) in the region.

%%%%%%%%%%%%%%%%%%%%%%%%%%%%%%%%%%%%%%%%%%%%%%%%%%%%%%%%%%%%
\subsection{Cauchy's Integral Formula}
\label{subsec:cauchy-integral-formula-summary}
%%%%%%%%%%%%%%%%%%%%%%%%%%%%%%%%%%%%%%%%%%%%%%%%%%%%%%%%%%%%

If \(f\) is holomorphic on and inside a positively oriented simple closed
contour \(C\), then for \(a\) inside \(C\),

\[
\boxed{
f(a)
=
\frac{
1
}{
2\pi i
}
\oint_C
\frac{
f(z)
}{
z-a
}
\,dz.
}
\]

%%%%%%%%%%%%%%%%%%%%%%%%%%%%%%%%%%%%%%%%%%%%%%%%%%%%%%%%%%%%
\subsection{Liouville's Theorem}
\label{subsec:liouville-summary}
%%%%%%%%%%%%%%%%%%%%%%%%%%%%%%%%%%%%%%%%%%%%%%%%%%%%%%%%%%%%

Every bounded entire function is constant.

This theorem provides a short proof of the Fundamental Theorem of Algebra.

%%%%%%%%%%%%%%%%%%%%%%%%%%%%%%%%%%%%%%%%%%%%%%%%%%%%%%%%%%%%
\subsection{Maximum Modulus Principle}
\label{subsec:maximum-modulus-summary}
%%%%%%%%%%%%%%%%%%%%%%%%%%%%%%%%%%%%%%%%%%%%%%%%%%%%%%%%%%%%

A nonconstant holomorphic function cannot attain a local maximum of its
modulus in the interior of a connected domain.

%%%%%%%%%%%%%%%%%%%%%%%%%%%%%%%%%%%%%%%%%%%%%%%%%%%%%%%%%%%%
\subsection{Residue Theorem}
\label{subsec:residue-theorem-summary}
%%%%%%%%%%%%%%%%%%%%%%%%%%%%%%%%%%%%%%%%%%%%%%%%%%%%%%%%%%%%

If \(f\) is meromorphic inside and on a suitable closed contour \(C\), then

\[
\boxed{
\oint_C
f(z)\,dz
=
2\pi i
\sum_k
\operatorname{Res}(f;z_k),
}
\]

where the sum is over poles inside \(C\).

%%%%%%%%%%%%%%%%%%%%%%%%%%%%%%%%%%%%%%%%%%%%%%%%%%%%%%%%%%%%
\subsection{Picard--Lindelöf Theorem}
\label{subsec:picard-lindelof-summary}
%%%%%%%%%%%%%%%%%%%%%%%%%%%%%%%%%%%%%%%%%%%%%%%%%%%%%%%%%%%%

For an initial-value problem

\[
y'=f(t,y),
\qquad
y(t_0)=y_0,
\]

appropriate continuity in \(t\) and local Lipschitz continuity in \(y\)
guarantee a locally unique solution.

%%%%%%%%%%%%%%%%%%%%%%%%%%%%%%%%%%%%%%%%%%%%%%%%%%%%%%%%%%%%
\subsection{Existence and Uniqueness for Linear ODEs}
\label{subsec:linear-ode-existence-summary}
%%%%%%%%%%%%%%%%%%%%%%%%%%%%%%%%%%%%%%%%%%%%%%%%%%%%%%%%%%%%

For

\[
y'+p(t)y=q(t),
\]

if \(p,q\) are continuous on an interval containing \(t_0\), then the initial
condition

\[
y(t_0)=y_0
\]

determines a unique solution throughout that interval.

%%%%%%%%%%%%%%%%%%%%%%%%%%%%%%%%%%%%%%%%%%%%%%%%%%%%%%%%%%%%
\subsection{Law of Large Numbers}
\label{subsec:law-large-numbers-summary}
%%%%%%%%%%%%%%%%%%%%%%%%%%%%%%%%%%%%%%%%%%%%%%%%%%%%%%%%%%%%

For independent identically distributed random variables under appropriate
moment assumptions, the sample mean approaches the population mean.

A standard weak-law statement is

\[
\boxed{
\bar X_n
\xrightarrow{P}
\mu.
}
\]

A strong-law version gives almost-sure convergence under suitable conditions.

%%%%%%%%%%%%%%%%%%%%%%%%%%%%%%%%%%%%%%%%%%%%%%%%%%%%%%%%%%%%
\subsection{Central Limit Theorem}
\label{subsec:central-limit-theorem-summary}
%%%%%%%%%%%%%%%%%%%%%%%%%%%%%%%%%%%%%%%%%%%%%%%%%%%%%%%%%%%%

For independent identically distributed random variables with finite mean
\(\mu\) and positive finite variance \(\sigma^2\),

\[
\boxed{
\frac{
\sum_{i=1}^{n}X_i-n\mu
}{
\sigma\sqrt n
}
\Longrightarrow
N(0,1).
}
\]

Equivalently,

\[
\boxed{
\frac{
\sqrt n(\bar X_n-\mu)
}{
\sigma
}
\Longrightarrow
N(0,1).
}
\]

%%%%%%%%%%%%%%%%%%%%%%%%%%%%%%%%%%%%%%%%%%%%%%%%%%%%%%%%%%%%
\subsection{Bayes' Theorem}
\label{subsec:bayes-theorem-summary}
%%%%%%%%%%%%%%%%%%%%%%%%%%%%%%%%%%%%%%%%%%%%%%%%%%%%%%%%%%%%

If \(\Pr(B)>0\),

\[
\boxed{
\Pr(A\mid B)
=
\frac{
\Pr(B\mid A)\Pr(A)
}{
\Pr(B)
}.
}
\]

For a partition \(A_1,\ldots,A_n\),

\[
\boxed{
\Pr(A_i\mid B)
=
\frac{
\Pr(B\mid A_i)\Pr(A_i)
}{
\sum_j
\Pr(B\mid A_j)\Pr(A_j)
}.
}
\]

%%%%%%%%%%%%%%%%%%%%%%%%%%%%%%%%%%%%%%%%%%%%%%%%%%%%%%%%%%%%
\subsection{Law of Total Probability}
\label{subsec:total-probability-summary}
%%%%%%%%%%%%%%%%%%%%%%%%%%%%%%%%%%%%%%%%%%%%%%%%%%%%%%%%%%%%

If \(A_1,\ldots,A_n\) form a partition, then

\[
\boxed{
\Pr(B)
=
\sum_{i=1}^{n}
\Pr(B\mid A_i)
\Pr(A_i).
}
\]

%%%%%%%%%%%%%%%%%%%%%%%%%%%%%%%%%%%%%%%%%%%%%%%%%%%%%%%%%%%%
\subsection{Law of Total Expectation}
\label{subsec:total-expectation-summary}
%%%%%%%%%%%%%%%%%%%%%%%%%%%%%%%%%%%%%%%%%%%%%%%%%%%%%%%%%%%%

For suitable random variables,

\[
\boxed{
\mathbb{E}[X]
=
\mathbb{E}
[
\mathbb{E}[X\mid Y]
].
}
\]

This is also called the tower property.

%%%%%%%%%%%%%%%%%%%%%%%%%%%%%%%%%%%%%%%%%%%%%%%%%%%%%%%%%%%%
\subsection{Law of Total Variance}
\label{subsec:total-variance-summary}
%%%%%%%%%%%%%%%%%%%%%%%%%%%%%%%%%%%%%%%%%%%%%%%%%%%%%%%%%%%%

Under appropriate integrability,

\[
\boxed{
\operatorname{Var}(X)
=
\mathbb{E}
[
\operatorname{Var}(X\mid Y)
]
+
\operatorname{Var}
(
\mathbb{E}[X\mid Y]
).
}
\]

%%%%%%%%%%%%%%%%%%%%%%%%%%%%%%%%%%%%%%%%%%%%%%%%%%%%%%%%%%%%
\subsection{Markov's Inequality}
\label{subsec:markov-inequality-summary}
%%%%%%%%%%%%%%%%%%%%%%%%%%%%%%%%%%%%%%%%%%%%%%%%%%%%%%%%%%%%

For nonnegative random variable \(X\) and \(a>0\),

\[
\boxed{
\Pr(X\geq a)
\leq
\frac{
\mathbb{E}[X]
}{
a
}.
}
\]

%%%%%%%%%%%%%%%%%%%%%%%%%%%%%%%%%%%%%%%%%%%%%%%%%%%%%%%%%%%%
\subsection{Chebyshev's Inequality}
\label{subsec:chebyshev-inequality-summary}
%%%%%%%%%%%%%%%%%%%%%%%%%%%%%%%%%%%%%%%%%%%%%%%%%%%%%%%%%%%%

For finite mean \(\mu\) and variance \(\sigma^2\),

\[
\boxed{
\Pr(
|X-\mu|
\geq
k\sigma
)
\leq
\frac1{k^2}.
}
\]

Equivalently,

\[
\boxed{
\Pr(
|X-\mu|
\geq
a
)
\leq
\frac{
\sigma^2
}{
a^2
}.
}
\]

%%%%%%%%%%%%%%%%%%%%%%%%%%%%%%%%%%%%%%%%%%%%%%%%%%%%%%%%%%%%
\subsection{Jensen's Inequality}
\label{subsec:jensen-inequality-summary}
%%%%%%%%%%%%%%%%%%%%%%%%%%%%%%%%%%%%%%%%%%%%%%%%%%%%%%%%%%%%

For convex \(\varphi\),

\[
\boxed{
\varphi(
\mathbb{E}[X]
)
\leq
\mathbb{E}
[
\varphi(X)
]
}
\]

under appropriate integrability assumptions.

For concave functions, the inequality reverses.

%%%%%%%%%%%%%%%%%%%%%%%%%%%%%%%%%%%%%%%%%%%%%%%%%%%%%%%%%%%%
\subsection{Cramér--Rao Lower Bound}
\label{subsec:cramer-rao-summary}
%%%%%%%%%%%%%%%%%%%%%%%%%%%%%%%%%%%%%%%%%%%%%%%%%%%%%%%%%%%%

Under standard regularity assumptions, an unbiased estimator \(\hat\theta\)
satisfies

\[
\boxed{
\operatorname{Var}(\hat\theta)
\geq
\frac{
1
}{
I(\theta)
},
}
\]

where \(I(\theta)\) is Fisher information for the relevant sample model.

%%%%%%%%%%%%%%%%%%%%%%%%%%%%%%%%%%%%%%%%%%%%%%%%%%%%%%%%%%%%
\subsection{Gauss--Markov Theorem}
\label{subsec:gauss-markov-summary}
%%%%%%%%%%%%%%%%%%%%%%%%%%%%%%%%%%%%%%%%%%%%%%%%%%%%%%%%%%%%

Under the classical linear regression assumptions, ordinary least squares is
the best linear unbiased estimator of the regression coefficients.

Here ``best'' means minimum variance within the class of linear unbiased
estimators.

%%%%%%%%%%%%%%%%%%%%%%%%%%%%%%%%%%%%%%%%%%%%%%%%%%%%%%%%%%%%
\subsection{Neyman--Pearson Lemma}
\label{subsec:neyman-pearson-summary}
%%%%%%%%%%%%%%%%%%%%%%%%%%%%%%%%%%%%%%%%%%%%%%%%%%%%%%%%%%%%

For testing a simple null hypothesis against a simple alternative, a
likelihood-ratio rejection region gives a most powerful test of a specified
size.

%%%%%%%%%%%%%%%%%%%%%%%%%%%%%%%%%%%%%%%%%%%%%%%%%%%%%%%%%%%%
\subsection{Max-Flow Min-Cut Theorem}
\label{subsec:max-flow-min-cut-summary}
%%%%%%%%%%%%%%%%%%%%%%%%%%%%%%%%%%%%%%%%%%%%%%%%%%%%%%%%%%%%

In a capacitated flow network,

\[
\boxed{
\text{maximum flow value}
=
\text{minimum cut capacity}.
}
\]

%%%%%%%%%%%%%%%%%%%%%%%%%%%%%%%%%%%%%%%%%%%%%%%%%%%%%%%%%%%%
\subsection{Hall's Marriage Theorem}
\label{subsec:halls-marriage-summary}
%%%%%%%%%%%%%%%%%%%%%%%%%%%%%%%%%%%%%%%%%%%%%%%%%%%%%%%%%%%%

A bipartite graph with parts \(X,Y\) has a matching that covers every vertex
of \(X\) if and only if

\[
\boxed{
|N(S)|
\geq
|S|
}
\]

for every subset \(S\subseteq X\).

%%%%%%%%%%%%%%%%%%%%%%%%%%%%%%%%%%%%%%%%%%%%%%%%%%%%%%%%%%%%
\subsection{Euler's Formula for Planar Graphs}
\label{subsec:euler-planar-summary}
%%%%%%%%%%%%%%%%%%%%%%%%%%%%%%%%%%%%%%%%%%%%%%%%%%%%%%%%%%%%

For a connected planar graph embedded in the plane,

\[
\boxed{
V-E+F=2.
}
\]

For multiple connected components, the corresponding form becomes

\[
\boxed{
V-E+F
=
1+C,
}
\]

where \(C\) is the number of connected components.

%%%%%%%%%%%%%%%%%%%%%%%%%%%%%%%%%%%%%%%%%%%%%%%%%%%%%%%%%%%%
\subsection{Euler's Theorem for Eulerian Circuits}
\label{subsec:eulerian-circuit-theorem-summary}
%%%%%%%%%%%%%%%%%%%%%%%%%%%%%%%%%%%%%%%%%%%%%%%%%%%%%%%%%%%%

A finite connected graph has an Eulerian circuit if and only if every vertex
has even degree.

A connected graph has an Eulerian trail with distinct endpoints if and only
if exactly two vertices have odd degree.

%%%%%%%%%%%%%%%%%%%%%%%%%%%%%%%%%%%%%%%%%%%%%%%%%%%%%%%%%%%%
\subsection{Handshaking Lemma}
\label{subsec:handshaking-lemma-summary}
%%%%%%%%%%%%%%%%%%%%%%%%%%%%%%%%%%%%%%%%%%%%%%%%%%%%%%%%%%%%

For any finite undirected graph,

\[
\boxed{
\sum_{v\in V}
\deg(v)
=
2|E|.
}
\]

Therefore, the number of vertices of odd degree is even.

%%%%%%%%%%%%%%%%%%%%%%%%%%%%%%%%%%%%%%%%%%%%%%%%%%%%%%%%%%%%
\subsection{Pigeonhole Principle}
\label{subsec:pigeonhole-principle-summary}
%%%%%%%%%%%%%%%%%%%%%%%%%%%%%%%%%%%%%%%%%%%%%%%%%%%%%%%%%%%%

If more than \(n\) objects are placed into \(n\) boxes, then at least one box
contains at least two objects.

A generalized form says that placing \(N\) objects into \(k\) boxes forces
some box to contain at least

\[
\boxed{
\left\lceil
\frac{N}{k}
\right\rceil
}
\]

objects.

%%%%%%%%%%%%%%%%%%%%%%%%%%%%%%%%%%%%%%%%%%%%%%%%%%%%%%%%%%%%
\subsection{Inclusion--Exclusion Principle}
\label{subsec:inclusion-exclusion-summary}
%%%%%%%%%%%%%%%%%%%%%%%%%%%%%%%%%%%%%%%%%%%%%%%%%%%%%%%%%%%%

For two finite sets,

\[
\boxed{
|A\cup B|
=
|A|+|B|-|A\cap B|.
}
\]

For three sets,

\[
\boxed{
|A\cup B\cup C|
=
|A|+|B|+|C|
-
|A\cap B|
-
|A\cap C|
-
|B\cap C|
+
|A\cap B\cap C|.
}
\]

The general principle alternates intersection terms by order.

%%%%%%%%%%%%%%%%%%%%%%%%%%%%%%%%%%%%%%%%%%%%%%%%%%%%%%%%%%%%
\subsection{Burnside's Lemma}
\label{subsec:burnside-lemma-summary}
%%%%%%%%%%%%%%%%%%%%%%%%%%%%%%%%%%%%%%%%%%%%%%%%%%%%%%%%%%%%

If finite group \(G\) acts on finite set \(X\), the number of orbits is

\[
\boxed{
|X/G|
=
\frac1{|G|}
\sum_{g\in G}
|\operatorname{Fix}(g)|.
}
\]

%%%%%%%%%%%%%%%%%%%%%%%%%%%%%%%%%%%%%%%%%%%%%%%%%%%%%%%%%%%%
\subsection{Euler--Lagrange Equation}
\label{subsec:euler-lagrange-summary}
%%%%%%%%%%%%%%%%%%%%%%%%%%%%%%%%%%%%%%%%%%%%%%%%%%%%%%%%%%%%

For a functional

\[
J[y]
=
\int
L(x,y,y')
\,dx,
\]

a sufficiently smooth extremizer satisfies

\[
\boxed{
\frac{\partial L}{\partial y}
-
\frac{d}{dx}
\left(
\frac{\partial L}{\partial y'}
\right)
=
0.
}
\]

%%%%%%%%%%%%%%%%%%%%%%%%%%%%%%%%%%%%%%%%%%%%%%%%%%%%%%%%%%%%
\subsection{Lagrange Multiplier Condition}
\label{subsec:lagrange-multiplier-summary}
%%%%%%%%%%%%%%%%%%%%%%%%%%%%%%%%%%%%%%%%%%%%%%%%%%%%%%%%%%%%

For an equality-constrained optimization problem

\[
\min f(x)
\]

subject to

\[
g(x)=0,
\]

a regular local optimum satisfies

\[
\boxed{
\nabla f
=
\lambda\nabla g.
}
\]

With several constraints,

\[
\boxed{
\nabla f
=
\sum_i
\lambda_i
\nabla g_i.
}
\]

%%%%%%%%%%%%%%%%%%%%%%%%%%%%%%%%%%%%%%%%%%%%%%%%%%%%%%%%%%%%
\subsection{Karush--Kuhn--Tucker Conditions}
\label{subsec:kkt-summary}
%%%%%%%%%%%%%%%%%%%%%%%%%%%%%%%%%%%%%%%%%%%%%%%%%%%%%%%%%%%%

For differentiable constrained optimization, the KKT conditions combine:

\[
\boxed{
\text{stationarity},
}
\]

\[
\boxed{
\text{primal feasibility},
}
\]

\[
\boxed{
\text{dual feasibility},
}
\]

and

\[
\boxed{
\text{complementary slackness}.
}
\]

Under suitable convexity and constraint-qualification assumptions, these
conditions characterize optimality.

%%%%%%%%%%%%%%%%%%%%%%%%%%%%%%%%%%%%%%%%%%%%%%%%%%%%%%%%%%%%
\subsection{Weak Duality}
\label{subsec:weak-duality-summary}
%%%%%%%%%%%%%%%%%%%%%%%%%%%%%%%%%%%%%%%%%%%%%%%%%%%%%%%%%%%%

For a minimization primal problem and its dual, every dual feasible objective
value provides a lower bound on every primal feasible objective value.

Thus,

\[
\boxed{
d^\star
\leq
p^\star.
}
\]

%%%%%%%%%%%%%%%%%%%%%%%%%%%%%%%%%%%%%%%%%%%%%%%%%%%%%%%%%%%%
\subsection{Strong Duality}
\label{subsec:strong-duality-summary}
%%%%%%%%%%%%%%%%%%%%%%%%%%%%%%%%%%%%%%%%%%%%%%%%%%%%%%%%%%%%

Under suitable conditions, the optimal primal and dual values coincide:

\[
\boxed{
d^\star
=
p^\star.
}
\]

The precise hypotheses depend on the optimization framework.

%%%%%%%%%%%%%%%%%%%%%%%%%%%%%%%%%%%%%%%%%%%%%%%%%%%%%%%%%%%%
\subsection{Quick Theorem Classification Table}
\label{subsec:quick-theorem-classification}
%%%%%%%%%%%%%%%%%%%%%%%%%%%%%%%%%%%%%%%%%%%%%%%%%%%%%%%%%%%%

\begin{center}
\begin{tabular}{c|c}
\textbf{Area}
&
\textbf{Representative Theorems}
\\ \hline

Number Theory
&
Fundamental Theorem of Arithmetic, CRT, Fermat, Euler
\\

Linear Algebra
&
Rank--Nullity, Spectral Theorem, SVD
\\

Abstract Algebra
&
Lagrange, Sylow, Isomorphism Theorems
\\

Geometry
&
Pythagorean, Law of Sines, Law of Cosines
\\

Topology
&
Jordan Curve, Brouwer, Heine--Borel
\\

Calculus
&
IVT, MVT, Taylor, FTC
\\

Analysis
&
DCT, MCT, Hölder, Minkowski
\\

Complex Analysis
&
Cauchy, Liouville, Residue Theorem
\\

Probability
&
LLN, CLT, Bayes, Jensen
\\

Statistics
&
Gauss--Markov, Neyman--Pearson, Cramér--Rao
\\

Combinatorics
&
Pigeonhole, Inclusion--Exclusion, Burnside
\\

Graph Theory
&
Hall, Max-Flow Min-Cut, Euler Formula
\\

Optimization
&
Lagrange Multipliers, KKT, Duality
\end{tabular}
\end{center}

%%%%%%%%%%%%%%%%%%%%%%%%%%%%%%%%%%%%%%%%%%%%%%%%%%%%%%%%%%%%
\subsection{Common Errors}
\label{subsec:common-theorem-errors}
%%%%%%%%%%%%%%%%%%%%%%%%%%%%%%%%%%%%%%%%%%%%%%%%%%%%%%%%%%%%

\subsubsection{Using a Theorem Without Checking Hypotheses}

A correct theorem can produce an invalid conclusion when its assumptions are
not satisfied.

%%%%%%%%%%%%%%%%%%%%%%%%%%%%%%%%%%%%%%%%%%%%%%%%%%%%%%%%%%%%
\subsubsection{Confusing a Theorem With Its Converse}

From

\[
P\Rightarrow Q
\]

one cannot automatically conclude

\[
Q\Rightarrow P.
\]

The converse requires a separate theorem or proof.

%%%%%%%%%%%%%%%%%%%%%%%%%%%%%%%%%%%%%%%%%%%%%%%%%%%%%%%%%%%%
\subsubsection{Confusing Necessary and Sufficient Conditions}

If \(P\) is sufficient for \(Q\), then

\[
P\Rightarrow Q.
\]

If \(P\) is necessary for \(Q\), then

\[
Q\Rightarrow P.
\]

%%%%%%%%%%%%%%%%%%%%%%%%%%%%%%%%%%%%%%%%%%%%%%%%%%%%%%%%%%%%
\subsubsection{Ignoring Dimensional Assumptions}

Some results depend strongly on finite dimensionality.

For example, in \(\mathbb{R}^n\),

\[
\boxed{
\text{closed and bounded}
\iff
\text{compact},
}
\]

but this equivalence generally fails in infinite-dimensional normed spaces.

%%%%%%%%%%%%%%%%%%%%%%%%%%%%%%%%%%%%%%%%%%%%%%%%%%%%%%%%%%%%
\subsubsection{Ignoring Regularity Conditions}

Differentiability, continuity, measurability, integrability, compactness,
convexity, and completeness hypotheses are often mathematically essential.

%%%%%%%%%%%%%%%%%%%%%%%%%%%%%%%%%%%%%%%%%%%%%%%%%%%%%%%%%%%%
\subsection{Exercises}
\label{subsec:common-mathematical-theorems-exercises}
%%%%%%%%%%%%%%%%%%%%%%%%%%%%%%%%%%%%%%%%%%%%%%%%%%%%%%%%%%%%

\begin{exercise}[1]
State the Fundamental Theorem of Arithmetic.
\end{exercise}

\begin{exercise}[1]
State Bézout's identity.
\end{exercise}

\begin{exercise}[1]
State the Fundamental Theorem of Algebra.
\end{exercise}

\begin{exercise}[1]
State the Rank--Nullity Theorem.
\end{exercise}

\begin{exercise}[1]
State the Mean Value Theorem.
\end{exercise}

\begin{exercise}[1]
State the Fundamental Theorem of Calculus.
\end{exercise}

\begin{exercise}[1]
State the Central Limit Theorem in standardized form.
\end{exercise}

\begin{exercise}[1]
State Bayes' theorem.
\end{exercise}

\begin{exercise}[2]
Explain why continuity is required in the Intermediate Value Theorem.
\end{exercise}

\begin{exercise}[2]
Explain the difference between Brouwer's and Banach's fixed-point theorems.
\end{exercise}

\begin{exercise}[2]
Give an example where Heine--Borel applies.
\end{exercise}

\begin{exercise}[2]
Use the Pigeonhole Principle to prove that among \(13\) people, at least two
have birthdays in the same month.
\end{exercise}

\begin{exercise}[2]
Use Euler's planar graph formula on a simple planar graph.
\end{exercise}

\begin{exercise}[2]
Use Markov's inequality for a nonnegative random variable.
\end{exercise}

\begin{exercise}[2]
Use Chebyshev's inequality to bound a deviation probability.
\end{exercise}

\begin{exercise}[2]
Use Jensen's inequality with a convex function.
\end{exercise}

\begin{exercise}[3]
Show how Fermat's Little Theorem follows from Euler's theorem when the
modulus is prime.
\end{exercise}

\begin{exercise}[3]
Use the Spectral Theorem to diagonalize a real symmetric matrix.
\end{exercise}

\begin{exercise}[3]
Use the Inverse Function Theorem to justify local invertibility.
\end{exercise}

\begin{exercise}[3]
Use the Implicit Function Theorem to compute \(dy/dx\).
\end{exercise}

\begin{exercise}[3]
Use Green's theorem to convert a line integral into a double integral.
\end{exercise}

\begin{exercise}[3]
Use the Divergence Theorem to compute flux.
\end{exercise}

\begin{exercise}[3]
Use Stokes' theorem to relate circulation and curl.
\end{exercise}

\begin{exercise}[3]
Use the Max-Flow Min-Cut Theorem on a small network.
\end{exercise}

\begin{exercise}[3]
Use Hall's theorem to determine whether a perfect matching exists.
\end{exercise}

\begin{exercise}[3]
Use the KKT conditions on a simple convex optimization problem.
\end{exercise}

\begin{exercise}[4]
Compare the Monotone Convergence Theorem with the Dominated Convergence
Theorem.
\end{exercise}

\begin{exercise}[4]
Explain how Liouville's theorem implies the Fundamental Theorem of Algebra.
\end{exercise}

\begin{exercise}[4]
Compare weak and strong duality.
\end{exercise}

%%%%%%%%%%%%%%%%%%%%%%%%%%%%%%%%%%%%%%%%%%%%%%%%%%%%%%%%%%%%
\subsection{Conceptual Summary}
\label{subsec:common-theorems-summary}
%%%%%%%%%%%%%%%%%%%%%%%%%%%%%%%%%%%%%%%%%%%%%%%%%%%%%%%%%%%%

Major mathematical theorems connect broad areas of mathematics.

Algebraic structure is governed by results such as

\[
\boxed{
\text{Fundamental Theorem of Arithmetic},
}
\]

\[
\boxed{
\text{Rank--Nullity},
}
\]

and

\[
\boxed{
\text{Spectral Theorem}.
}
\]

Calculus and analysis rely on

\[
\boxed{
\text{Intermediate Value Theorem},
}
\]

\[
\boxed{
\text{Mean Value Theorem},
}
\]

\[
\boxed{
\text{Fundamental Theorem of Calculus},
}
\]

and convergence theorems.

Probability and statistics rely on

\[
\boxed{
\text{Law of Large Numbers},
\quad
\text{Central Limit Theorem},
\quad
\text{Bayes' Theorem}.
}
\]

Optimization relies on

\[
\boxed{
\text{Lagrange multipliers},
\quad
\text{KKT conditions},
\quad
\text{duality}.
}
\]

The central reference principle is

\[
\boxed{
\text{a theorem is valid only when its hypotheses are satisfied}.
}
\]

%%%%%%%%%%%%%%%%%%%%%%%%%%%%%%%%%%%%%%%%%%%%%%%%%%%%%%%%%%%%
\subsection{Section Connections}
\label{subsec:common-mathematical-theorems-connections}
%%%%%%%%%%%%%%%%%%%%%%%%%%%%%%%%%%%%%%%%%%%%%%%%%%%%%%%%%%%%

\chapterconnection{
Common Mathematical Theorems provides a cross-disciplinary index of major
results developed throughout the textbook. It connects foundational results
from Algebra, Number Theory, Geometry, Topology, Analysis, Probability,
Statistics, Differential Equations, Discrete Mathematics, and Optimization.
The next reference section, Mathematical Terminology, organizes frequently
used mathematical vocabulary by concept and subject area before the
alphabetical Glossary in Section~14.17.
}

%%%%%%%%%%%%%%%%%%%%%%%%%%%%%%%%%%%%%%%%%%%%%%%%%%%%%%%%%%%%
% SECTION SOURCE TRACKING
%%%%%%%%%%%%%%%%%%%%%%%%%%%%%%%%%%%%%%%%%%%%%%%%%%%%%%%%%%%%

%------------------------------------------------------------
% 14.15 Common Mathematical Theorems
%------------------------------------------------------------
%
% Source basis:
%   - Standard foundational theorems from undergraduate and
%     graduate mathematics.
%
% Used for:
%   - arithmetic and number theory;
%   - algebra and linear algebra;
%   - abstract algebra;
%   - geometry and topology;
%   - calculus and real analysis;
%   - functional analysis;
%   - complex analysis;
%   - differential equations;
%   - probability and statistics;
%   - combinatorics and graph theory;
%   - optimization and variational methods.
%
% Notes:
%   - Statements are intentionally compact.
%   - Full hypotheses, proofs, and applications remain in their
%     canonical subject chapters.
%   - Mathematical Terminology follows in Section 14.16.
%
%%%%%%%%%%%%%%%%%%%%%%%%%%%%%%%%%%%%%%%%%%%%%%%%%%%%%%%%%%%%